% !TeX root = ../概率论与数理统计笔记.tex
% 由主文档通过 \input 载入；请编译主文档。
\setcounter{chapter}{12} % 保留原书章号。
\chapter{区间估计与置信区域}
\label{chap:ps-intervals}

本章用区间或区域表达估计的不确定性，核心是重复抽样意义下的覆盖性质。
先修内容为抽样分布、点估计与中心极限定理；学习目标是区分精确区间与渐近区间，并正确解释置信水平。

\section{置信区间的定义与枢轴量}
\label{sec:ps-confidence-definition}
令 $0<\alpha<1$，定义\emph{置信区间}（confidence interval）$[L(\mathbf X),U(\mathbf X)]$ 的覆盖概率。
对精确的至少 $1-\alpha$ 水平，要求对每个 $\theta\in\Theta$ 有
\begin{equation}
  \mathbb P_\theta\{L(\mathbf X)\leq\theta\leq U(\mathbf X)\}\geq1-\alpha.
  \label{eq:ps-confidence-coverage}
\end{equation}
介绍\emph{枢轴量}（pivotal quantity）的构造与反解，区分随机区间和观测后的固定区间。
频率学派框架中的参数固定，不把置信水平直接解释为观测后参数落入该固定区间的概率。

\section{单个正态总体的参数区间}
\label{sec:ps-normal-intervals}
分别推导方差已知或未知时均值的双侧和单侧区间，以及正态总体方差的卡方区间。
明确正态性、独立同分布、样本量与自由度条件，统一分位数采用左侧累计概率的约定。
比较标准误、置信水平和区间长度之间的关系。

\section{两个总体的比较与配对数据}
\label{sec:ps-two-sample-intervals}
讨论独立正态样本的均值差区间，区分方差已知、未知但相等、未知且不等的情形。
给出合并方差的精确 $t$ 方法与韦尔奇近似（Welch approximation）的条件及区别。
配对样本先构造对内差值，再对独立的配对单位进行推断；方差比区间另需两个正态样本独立。

\section{比例、大样本区间与样本量规划}
\label{sec:ps-asymptotic-intervals}
用渐近正态性与标准误构造大样本区间，并以二项比例为例比较沃尔德区间（Wald interval）、威尔逊区间（Wilson interval）和精确反演区间。
讨论小样本、接近边界的参数与离散分布对覆盖率的影响；近似区间不承诺式\eqref{eq:ps-confidence-coverage}的有限样本覆盖。
根据目标误差和置信水平规划样本量，明确使用已知或预估方差的前提。

\section{例题、习题与学习检查}
规划正态均值与方差区间、两种测量方法的配对差区间、二项比例在小样本下的区间比较。
习题要求从枢轴量的不等式完整反解端点，检查单侧与双侧分位数。
置信区域作为向量参数的自然推广介绍；预测区间与均值区间的区别在第\ref{chap:ps-regression}章进一步展开。






